% !TEX root = ../main.tex

% 预习报告	
\setcounter{section}{0}
\section{锁相放大器工作原理、基本参数与基本操作 \\ Preview Report}


%---------------------------------------------------------------------
% 实验目的
\subsection{Purpose}
\begin{enumerate}
	\item 了解锁相放大器工作原理和特点，理解信号、噪声、信噪比等概念。
	\item 掌握锁相放大器基本参数含义及锁相放大器的基本操作，学会合理选择或调节参数（频率、相位、灵敏度、时间常数、陡降）；复习示波器的使用；
	\item 掌握用锁相放大器检测弱信号方法，通过与示波器比较其检测能力了解其技术优势。
\end{enumerate}


%---------------------------------------------------------------------
% 仪器用具
\subsection{Instruments \& Equipment}
\begin{table}[htbp]
	\centering
	\renewcommand\arraystretch{1.6}
	% \setlength{\tabcolsep}{10mm}
	\begin{tabular}{|p{0.08\textwidth}|p{0.16\textwidth}|p{0.08\textwidth}|p{0.6\textwidth}|}
		\hline
		Number & Name & Quantity & Main parameters (model, measurement range, measurement accuracy, etc.) \\
		\hline
		1 & 锁相放大器 & 1 & OE1022（系列） \\
		\hline
		2 & 配套教学实验箱 & 1 & \\
		\hline
		3 & 示波器 & 1 & RIGOL DS2202A \\
		\hline
		4 & 信号发生器 & 1 & RIGOL DG4162 \\
		\hline
		5 & BNC-BNC信号线 & 若干 & \\
		\hline
	\end{tabular}
\end{table}


%---------------------------------------------------------------------
% 原理概述
\subsection{Principle}
（由于预习报告不做要求，因此这里只简单总结原理中我觉得重要的部分，帮助理解实验环节。）
\begin{enumerate}
	\item 信号与噪声
	
	\textbf{概述}
	
	信号是随时间、空间或其他独立变量变化的物理量，用于携带信息。在数学上，信号通常表示为一个函数，例如连续时间信号 \( x(t) \) 和离散时间信号 \( x[n] \)。根据时间特性，信号可以分为连续信号和离散信号；根据幅度特性，信号又可分为连续幅度信号和离散幅度信号。常见的信号如正弦信号和指数信号，其中正弦信号的数学形式为 \( x(t) = A \sin(2\pi f t + \phi) \)，其中 \( A \) 为振幅，\( f \) 为频率，\( \phi \) 为相位；指数信号则表示为 \( x(t) = A e^{\sigma t} \)，其中 \( \sigma \) 是增长或衰减系数。
	
	噪声是干扰信号传输和处理的非期望或随机信号，通常会降低信号的质量和可用性。噪声来源广泛，包括热噪声（由电子元件的热运动产生）、散粒噪声（电流的离散性引起）以及环境噪声（外部环境干扰）。噪声可以通过不同的数学模型描述，常见的有白噪声和高斯噪声。白噪声是功率谱密度恒定的随机过程，通常表示为 \( n(t) \)；而高斯噪声则是服从高斯分布的随机过程，其概率密度函数为 \( p(n) = \frac{1}{\sqrt{2\pi \sigma^2}} \exp\left( -\frac{(n - \mu)^2}{2\sigma^2} \right) \)，其中 \( \mu \) 为均值，\( \sigma^2 \) 为方差。
	
	在实际系统中，接收到的信号通常是原始信号与噪声的叠加，即 \( r(t) = s(t) + n(t) \)，其中 \( r(t) \) 是接收信号，\( s(t) \) 是原始信号，\( n(t) \) 是噪声。信号与噪声的关系常用信噪比（Signal-to-Noise Ratio, SNR）来衡量，信噪比是信号功率与噪声功率的比值，用于衡量信号质量。其公式为 \( \mathrm{SNR} = \frac{P_{\text{signal}}}{P_{\text{noise}}} \)，并常使用分贝形式表示：\( \mathrm{SNR}_{\mathrm{dB}} = 10 \log_{10} \left( \frac{P_{\text{signal}}}{P_{\text{noise}}} \right) \)。较高的 SNR 表示信号质量较好，噪声对信号的影响较小。
	
	噪声具有统计特性，可以通过均值和自相关函数来描述。噪声的功率谱密度（PSD）描述了噪声在频域内的分布，其定义为 \( S_n(f) = \int_{-\infty}^{\infty} R_n(\tau) e^{-j2\pi f \tau} d\tau \)，其中 \( R_n(\tau) \) 是自相关函数。对于白噪声，功率谱密度是恒定的，即 \( S_n(f) = N_0 / 2 \)，其中 \( N_0 \) 是噪声功率谱密度。
	
	噪声对系统性能有显著影响。例如，在通信系统中，噪声会导致误码率增加，对于二进制相移键控（BPSK），误码率公式为 \( \mathrm{BER} = Q\left( \sqrt{\frac{2E_b}{N_0}} \right) \)，其中 \( E_b \) 是每比特能量，\( N_0 \) 是噪声功率谱密度，\( Q \) 函数为 \( Q(x) = \frac{1}{\sqrt{2\pi}} \int_x^\infty e^{-t^2/2} dt \)。此外，根据香农定理，信道容量也受到噪声的限制，其最大信息传输速率为 \( C = B \log_2\left(1 + \frac{P}{N_0 B}\right) \)，其中 \( C \) 是信道容量（比特每秒），\( B \) 是带宽，\( P \) 是信号功率。
	
	噪声的类型和模型在系统分析中也有重要意义，加性白高斯噪声（AWGN）是一个常见的噪声模型，用于分析通信系统的性能。在 AWGN 模型中，接收信号表示为 \( r(t) = s(t) + n(t) \)，其中 \( n(t) \) 是 AWGN。此外，窄带噪声是通过带限系统产生的噪声，其统计特性更加复杂。
	
	为了抑制噪声对信号的影响，可以采用滤波器设计、采样与量化以及信号平均等方法。滤波器设计中，低通滤波器可以抑制高频噪声，高通滤波器用于抑制低频噪声，带通滤波器则用于提取特定频段的信号。通过增加信号的平均次数，也可以提高信噪比，其关系为 \( \mathrm{SNR}_{\mathrm{average}} = \sqrt{N} \times \mathrm{SNR}_{\mathrm{single}} \)，其中 \( N \) 是平均次数。
	
	在实际应用中，信号与噪声的处理非常关键，例如在移动通信、卫星通信和图像处理领域。通信系统中，噪声会影响通话质量和数据传输速率，而图像处理中，常见的噪声如高斯噪声、椒盐噪声会影响图像质量，需要通过滤波等去噪算法进行处理。此外，在生物医学信号处理中，如心电信号（ECG），噪声会影响信号的精确提取，因此通常通过信号预处理技术来提取有效信号。
	
	\textbf{微弱信号}
	
	“微弱信号”多指相对于噪声而言幅度很小的信号，特别是那些被噪声淹没的信号。只有在有效地抑制噪声的条件下放大微弱信号的幅度，才能提取出真实的信号。提高信噪比的途径有两条：一是通过物理屏蔽来抑制环境噪声（压低强度），二是通过信号处理滤除噪声（缩小积分范围）。
	
	\textbf{什么叫信噪比（SNR）及信噪改善比（SNIR ）？  }
	
	信噪比（Signal-to-Noise Ratio, SNR）是衡量信号质量的重要指标，表示有用信号功率与噪声功率的比值，用于评估信号在噪声环境中的清晰度和可辨识度。信噪比的数学定义为：
	
	\[
	\mathrm{SNR} = \frac{P_{\text{signal}}}{P_{\text{noise}}}
	\]
	
	其中 \( P_{\text{signal}} \) 是信号的平均功率，\( P_{\text{noise}} \) 是噪声的平均功率。为了便于比较，通常将信噪比转换为分贝（dB）表示：
	
	\[
	\mathrm{SNR}_{\mathrm{dB}} = 10 \log_{10} \left( \frac{P_{\text{signal}}}{P_{\text{noise}}} \right)
	\]
	
	信噪比的物理意义是：较高的 SNR 表示信号功率远大于噪声功率，信号质量高，易于检测和处理；而较低的 SNR 表示噪声功率接近或超过信号功率，信号质量较差，可能导致误码或失真。在通信系统中，接收到的信号通常是原始信号与噪声的叠加，即 \( r(t) = s(t) + n(t) \)，其中 \( r(t) \) 为接收信号，\( s(t) \) 为有用信号，\( n(t) \) 为噪声。通过计算信号和噪声的功率，可以得出系统的信噪比，从而评估通信质量。
	
	信噪改善比（Signal-to-Noise Improvement Ratio, SNIR）用于衡量某种处理（如滤波、放大、编码等）对信噪比的提升程度。它表示处理后信噪比相对于处理前信噪比的提高倍数，数学定义为：
	
	\[
	\mathrm{SNIR} = \frac{\mathrm{SNR}_{\text{out}}}{\mathrm{SNR}_{\text{in}}}
	\]
	
	或以分贝形式表示为：
	
	\[
	\mathrm{SNIR}_{\mathrm{dB}} = \mathrm{SNR}_{\text{out, dB}} - \mathrm{SNR}_{\text{in, dB}}
	\]
	
	其中 \( \mathrm{SNR}_{\text{in}} \) 是处理前的信噪比，\( \mathrm{SNR}_{\text{out}} \) 是处理后的信噪比。信噪改善比的物理意义是，如果 SNIR > 1（或 \( \mathrm{SNIR}_{\mathrm{dB}} > 0 \)），表示信噪比得到改善，信号质量提高；如果 SNIR = 1（或 \( \mathrm{SNIR}_{\mathrm{dB}} = 0 \)），表示信噪比无变化；而如果 SNIR < 1（或 \( \mathrm{SNIR}_{\mathrm{dB}} < 0 \)），则信噪比降低，信号质量下降。例如，通过某种处理方法降低噪声功率而保持信号功率不变，则信噪改善比可以表示为：
	
	\[
	\mathrm{SNIR} = \frac{P_{\text{noise, in}}}{P_{\text{noise, out}}}
	\]
	
	这意味着信噪改善比等于噪声功率减少的倍数。
	
	以一个具体的实例来说明信噪比和信噪改善比的计算。假设一个系统中，信号功率为 \( P_{\text{signal}} = 50 \, \text{mW} \)，噪声功率为 \( P_{\text{noise}} = 5 \, \text{mW} \)。则信噪比为：
	
	\[
	\mathrm{SNR} = \frac{50}{5} = 10
	\]
	
	转换为分贝为：
	
	\[
	\mathrm{SNR}_{\mathrm{dB}} = 10 \log_{10}(10) = 10 \, \text{dB}
	\]
	
	如果通过某种处理方法将噪声功率降低到 \( P_{\text{noise, out}} = 1 \, \text{mW} \)，信号功率不变，则新的信噪比为：
	
	\[
	\mathrm{SNR}_{\text{out}} = \frac{50}{1} = 50
	\]
	
	转换为分贝为：
	
	\[
	\mathrm{SNR}_{\text{out, dB}} = 10 \log_{10}(50) \approx 16.99 \, \text{dB}
	\]
	
	信噪改善比为：
	
	\[
	\mathrm{SNIR} = \frac{50}{10} = 5
	\]
	
	或以分贝表示为：
	
	\[
	\mathrm{SNIR}_{\mathrm{dB}} = 16.99 \, \text{dB} - 10 \, \text{dB} = 6.99 \, \text{dB}
	\]
	
	这表示经过处理，信噪比提高了约 7 dB。
	
	\item 锁相放大技术
	
	\textbf{背景}
	
	采用零差检测背景方法和低通滤波方法，这种仪器能在极强噪声环境中提取信号幅值和相位信息。待测单元信号是一个周期信号，用锁相放大器可以测量幅值和相对于参考信号的相位。较强的信号示波器也能测，但是对于小信号，就会淹没在噪声中，在这种情况下，锁相放大器就起着重要作用。
	
	\begin{figure}[htbp]
		\centering
		\includegraphics[width=0.7\textwidth]{lockin_01_signal_buried_2023_final.png}
		\caption{锁相放大器可测量待测信号相对于特定参考信号的幅值和相位（甚至是完全淹没在噪声之中的信号）（参考文献[3]）。}
		\label{fig:lockin_01_signal_buried_2023_final}
	\end{figure}
	
	\textbf{混频}
	
	将输入信号和参考信号混频，混频其实就是相乘，当被测信号和参考信号都是正弦波时，混频的结果就是相加和相减。如果两信号频率相同，则获得直流和两倍频的信号，通过低通滤波器将直流和倍频分离。
	
	\textbf{滤波器}
	
	在-3dB此处信号的能量衰减一半，信号的带宽和时间常数成反比，带宽太大，要考虑信通性的误差，倍频部分可能会泄露到输出信号。此外，更大的带宽意味着更多的噪声，和更低的信噪比。如果滤波器带宽太窄，就会限制滤波器的响应时间，使得测量变的缓慢。低通滤波器的形状可以通过阶数来决定，更高的阶数更趋近于矩形的滤波函数，能够有效阻断滤波器外的频率，但高阶滤波器需要更高的时间达到稳定，导致相位延迟。降低阶数的好处是使相位延迟降低，时间延迟更短，低通滤波器更适合于高速信号的测量，混频加低通滤波就是向明检测的原理或者解调。
	
	\begin{figure}[htbp]
		\centering
		\includegraphics[width=0.5\textwidth]{IMG_3588.png}
		\caption{低通滤波器基本原理}
		\label{fig:IMG_3588}
	\end{figure}
	
	\begin{figure}[htbp]
		\centering
		\includegraphics[width=0.5\textwidth]{lockin_06_cascaded_filters_2023_final.png}
		\caption{(a) 一阶 RC 滤波器及其传递函数公式。(b) 通过将多个 RC 滤波器级联可以针对更高频率获得更陡的滚降特性。最终的传递函数为各滤波器传递函数相乘。（来自参考文献[3]）}
		\label{fig:lockin_06_cascaded_filters_2023_final}
	\end{figure}
	
	\begin{figure}[htbp]
		\centering
		\includegraphics[width=0.7\textwidth]{lockin_07_filters_same_tau_2023_final.png}
		\caption{图中的蓝线为 RC 滤波器传递函数的波特图。 同时，还绘制了具有相同滤波时间常数的高阶滤波器 (n = 2, 4, 8) 的传递函数，从图中可以看出，高阶滤波器的信号带宽 −3dB 相对要低得多。(c) 时域中对应的阶跃响应函数。通过级联多个滤波器，可以显著增加达到同一精度的稳定时间。这与相位延迟的增大有关。级联 RC 或积分滤波器的另一个优（甚至是完全淹没在噪声之中的信号）（来自参考文献[3]。}
		\label{fig:lockin_07_filters_same_tau_2023_final}
	\end{figure}
	
	\begin{figure}[htbp]
		\centering
		\includegraphics[width=0.7\textwidth]{lockin_08_filters_same_cutoff_2023_final.png}
		\caption{(a) 高阶滤波器针对更高频率的滚降更陡。(b) 高阶滤波器的相位延迟更大，会对反馈应用造成不良影响。(c) 时域内的阶跃响应函数，以一阶滤波器时间常数为单位。尽管低阶滤波器在初始时对输入信号变化的响应更为迅速，但这个优势会随时间的推移而逐渐减小，在某一时间点，高阶滤波器甚至会超越低阶滤波器，如图所示。（来自参考文献[3]）}
		\label{fig:lockin_08_filters_same_cutoff_2023_final}
	\end{figure}
	
	\textbf{双相解调}
	
	$$
	\begin{aligned}
		& X=\operatorname{Re}(Z)=\left\langle V_s(t) \cos \left(\omega_s t\right)\right\rangle=R \cos \Theta, \\
		& Y=\operatorname{Im}(Z)=-\left\langle V_s(t) \sin \left(\omega_s t\right)\right\rangle=R \sin \Theta .
	\end{aligned}
	$$
	
	$$
	\begin{aligned}
		& R=\sqrt{X^2+Y^2} \\
		& \Theta=\operatorname{atan} 2(Y, X)
	\end{aligned}
	$$
	
	\begin{figure}[htbp]
		\centering
		\includegraphics[width=0.6\textwidth]{lockin_02_dual_phase_2023_final.png}
		\caption{锁相测量过程的示意图（来自参考文献[3]）}
		\label{fig:lockin_02_dual_phase_2023_final}
	\end{figure}
	
	\textbf{利用傅里叶去理解解调非正弦信号}
	
	每一个周期信号等于无数个cos和sin，经过低通滤波器，如前面所说，信号被称为同相分量，标记为X， 与之对应的是，余弦信号解调出两个傅里叶分量中的一个，当我们用正弦信号解调时，即对参考信号相位延迟90度，我们得到另一个独立的线性分量，通常标记为Y，称为正交分量，当我们有了X和Y，锁相放大器就能计算出幅值和相位 ，把笛卡尔坐标系转化为极坐标系就可以了。
	
	
\end{enumerate}


%---------------------------------------------------------------------
% 实验前思考题
\subsection{Thinking Before Experiment}

\subsubsection{Reflection Questions of Principle}

% 思考题
\begin{question}
	市频 50Hz 干扰通常通过电源耦合，影响仪器的测量结果；对于 997Hz 的待测信号，50Hz 干扰是噪声吗？对锁相放大器的测量会有影响吗？
\end{question}
市频 50Hz 干扰主要通过电源线或其他电磁耦合路径引入到测量系统中，表现为一种低频的正弦信号。由于 50Hz 是非常常见的环境噪声源，许多实验仪器在测量过程中都会受到此频率的干扰。在实验中，50Hz 干扰可以被视为一种不相关的噪声。它与 997Hz 的待测信号频率明显不同，因此可以认为是一种外部噪声源。锁相放大器的一个主要功能就是从噪声背景中提取出目标信号，所以关键问题是 50Hz 的干扰是否会影响锁相放大器的测量。

待测信号的频率为 997Hz，而 50Hz 干扰频率与待测信号相差很大。理论上，锁相放大器会根据设置的参考频率（997Hz）来进行锁定，因此 50Hz 干扰会被滤除，并不会直接影响测量结果。

虽然锁相放大器能够在频率上选择性滤除不相关的噪声，但 50Hz 干扰仍有可能产生一些间接的影响。50Hz 干扰可能会产生 100Hz、150Hz 等谐波成分，这些频率成分接近或落在待测信号附近时，可能会引起较小的干扰，尤其是当谐波频率接近参考频率时，滤波器的效果会降低。50Hz 干扰可能引入低频的基线漂移，增加测量的不稳定性。如果干扰信号的幅度较大，可能会使得信号源的供电或参考信号不稳定，从而影响实验的重复性。

尽管锁相放大器能部分解决这个问题，但还是有一些额外的措施可以进一步降低50Hz干扰的影响。确保实验设备和电缆有良好的屏蔽，并做好接地工作，以减少电磁耦合带来的干扰。在锁相放大器的输入端加装陷波滤波器（50Hz陷波器），专门抑制 50Hz 的干扰信号。尽量远离大型电气设备或电源线，以减少市电干扰的耦合。

% 思考题
\begin{question}
	如何用锁相放大器检测到待测的直流信号或慢变信号？
\end{question}
锁相放大器的主要功能是从噪声背景中提取特定频率的信号。它通过将输入信号与参考信号（通常是一个正弦波）相乘，生成调制信号。然后，低通滤波器滤除与参考频率不同的成分，只保留同频率且相位一致的部分。这使锁相放大器在高噪声环境中非常有效，尤其是对于交流（AC）信号的提取。然而，锁相放大器的这种原理主要是为交流信号设计的，因此对于直流（DC）信号或缓慢变化的信号（低频信号），需要采取一些额外措施来适应其检测。

锁相放大器通常依赖于参考信号来选择性放大输入信号中的某个频率。如果待测信号是直流信号或频率非常低，它不会与常规交流参考信号匹配，难以直接应用常规的锁相放大技术。在直流信号情况下，如果直接通过锁相放大器来测量，由于没有与直流信号相对应的调制参考信号，锁相放大器会无法区分直流信号和直流偏置的噪声。因此，直接检测直流信号的有效性有限。

要检测直流或慢变信号，常用的方法是对信号进行调制，将其变为交流信号，然后使用锁相放大器检测调制后的信号。

斩波调制（Chopping modulation）是一种常见的方法，尤其在光学检测中。例如，假设要测量某个光电探测器的输出信号，该输出信号可能是缓慢变化的或接近直流。通过使用一个机械或电子斩波器，信号可以被调制为一个周期性的方波或正弦波。斩波器的频率通常设置在数百赫兹到千赫兹的范围内，这个频率同时作为锁相放大器的参考频率。经过调制后的信号从直流或慢变信号变成了一个与参考信号相匹配的交流信号，锁相放大器可以通过对调制频率进行锁定，检测出该调制后的信号。调制的优点是将低频或直流信号从低频噪声（如 1/f 噪声）中分离出来，显著提高信噪比。

另一种方法是引入一个外部调制器，使信号周期性地变化。例如，给待测的直流信号添加一个已知频率的小振幅交流调制。这个交流调制信号将与直流信号一起被输入到锁相放大器中，通过锁相放大器锁定该已知频率，可以提取出信号的幅度和相位信息。该方法广泛应用于许多领域，比如在激光光学系统中，通过给激光源施加调制，可以从背景噪声中提取目标信号。

对于慢变信号，虽然其频率较低，但仍是一个交流信号，可以通过适当的锁相放大器设置进行检测。锁相放大器的参考频率应与待测慢变信号的频率相匹配。由于信号变化缓慢，因此可以将参考频率设定得较低，甚至可以是毫赫兹级别。这样可以确保信号与参考信号的相位匹配，锁相放大器仍能有效检测到信号。对于慢变信号，锁相放大器的低通滤波器时间常数应适当增加。较长的时间常数可以减少高频噪声干扰，稳定慢变信号的检测。

检测直流或慢变信号时，由于信号频率低，系统中低频噪声的影响会显得更加明显。例如，1/f 噪声和温度漂移可能成为主要干扰源。因此，良好的屏蔽、接地和温控是保证实验精度的关键。在调制方法中，确保调制器和锁相放大器的同步性至关重要。如果两者不同步，信号提取可能会失败或产生较大误差。

% 思考题
\begin{question}
	如用斩波器调制直流信号（如光强），被斩制后的信号仍然包含有直流分量（即平均值不为零），但该直流分量随交流信号输入锁相放大器后不会被锁相放大器检测，请从数学推导上说明。
\end{question}
首先，直流信号（例如光强）经过斩波器调制，变成了一个周期性的交流信号。斩波器通常会将信号切换为一个方波或其他周期性波形。

设原始直流信号为：

\[
S_{\text{DC}}(t) = A
\]

其中，\( A \) 是直流信号的幅度。

斩波器的调制函数 \( M(t) \) 可以表示为一个方波，其值在 0 和 1 之间变化，频率为 \( f_{\text{chop}} \)：

\[
M(t) = \dfrac{1}{2} + \dfrac{2}{\pi} \sum_{n=1,3,5,\dots}^{\infty} \dfrac{1}{n} \sin(2\pi n f_{\text{chop}} t)
\]

因此，经过斩波器后的信号 \( S_{\text{chop}}(t) \) 为：

\[
S_{\text{chop}}(t) = S_{\text{DC}}(t) \cdot M(t) = A \cdot M(t)
\]

将 \( M(t) \) 展开为傅里叶级数后，发现它包含一个直流分量和一系列谐波分量。

斩波信号 \( S_{\text{chop}}(t) \) 可表示为：

\[
S_{\text{chop}}(t) = A \left( \dfrac{1}{2} + \dfrac{2}{\pi} \sum_{n=1,3,5,\dots}^{\infty} \dfrac{1}{n} \sin(2\pi n f_{\text{chop}} t) \right)
\]

其中，直流分量为\( \dfrac{A}{2} \)，交流分量包含频率为 \( f_{\text{chop}} \) 的基波和其奇数次谐波。

参考信号 \( R(t) \) 通常是一个与斩波器同步的正弦波：

\[
R(t) = \sin(2\pi f_{\text{ref}} t + \phi)
\]

由于斩波器的频率就是我们希望检测的信号频率，所以 \( f_{\text{ref}} = f_{\text{chop}} \)。

为了简化，设相位差 \( \phi = 0 \)：

\[
R(t) = \sin(2\pi f_{\text{chop}} t)
\]

将 \( S_{\text{chop}}(t) \) 与 \( R(t) \) 相乘，得到乘法器输出信号 \( S_{\text{mix}}(t) \)：

\[
S_{\text{mix}}(t) = S_{\text{chop}}(t) \cdot R(t) = A \cdot M(t) \cdot \sin(2\pi f_{\text{chop}} t)
\]

将 \( M(t) \) 展开并代入：

\[
S_{\text{mix}}(t) = A \left( \dfrac{1}{2} + \dfrac{2}{\pi} \sum_{n=1,3,5,\dots}^{\infty} \dfrac{1}{n} \sin(2\pi n f_{\text{chop}} t) \right) \sin(2\pi f_{\text{chop}} t)
\]

展开后：

\[
S_{\text{mix}}(t) = \dfrac{A}{2} \sin(2\pi f_{\text{chop}} t) + \dfrac{2A}{\pi} \sum_{n=1,3,5,\dots}^{\infty} \dfrac{1}{n} \sin(2\pi n f_{\text{chop}} t) \sin(2\pi f_{\text{chop}} t)
\]

利用三角恒等式：

\[
\sin A \sin B = \dfrac{1}{2} [\cos(A - B) - \cos(A + B)]
\]

对第二项进行展开：

\[
\sin(2\pi n f_{\text{chop}} t) \sin(2\pi f_{\text{chop}} t) = \dfrac{1}{2} \left[ \cos(2\pi (n - 1) f_{\text{chop}} t) - \cos(2\pi (n + 1) f_{\text{chop}} t) \right]
\]

因此，乘法器输出可表示为：

\[
\begin{aligned}
	S_{\text{mix}}(t) = & \dfrac{A}{2} \sin(2\pi f_{\text{chop}} t) \\
	& + \dfrac{2A}{\pi} \sum_{n=1,3,5,\dots}^{\infty} \dfrac{1}{n} \cdot \dfrac{1}{2} \left[ \cos(2\pi (n - 1) f_{\text{chop}} t) - \cos(2\pi (n + 1) f_{\text{chop}} t) \right] \\
	= & \dfrac{A}{2} \sin(2\pi f_{\text{chop}} t) + \dfrac{A}{\pi} \sum_{n=1,3,5,\dots}^{\infty} \dfrac{1}{n} \left[ \cos(2\pi (n - 1) f_{\text{chop}} t) - \cos(2\pi (n + 1) f_{\text{chop}} t) \right]
\end{aligned}
\]

低通滤波器的输入信号是 \( S_{\text{mix}}(t) \)，其中包含以下频率分量：第一项，\( \dfrac{A}{2} \sin(2\pi f_{\text{chop}} t) \) —— 频率为 \( f_{\text{chop}} \) 的交流分量。第二项，由谐波生成的一系列余弦函数，频率为 \( (n \pm 1) f_{\text{chop}} \)。

为了找到低通滤波器输出的直流分量，我们需要找出在 \( S_{\text{mix}}(t) \) 中频率为零的项。

观察第二项中的 \( \cos(2\pi (n - 1) f_{\text{chop}} t) \)，当 \( n = 1 \) 时：

\[
\cos(2\pi (1 - 1) f_{\text{chop}} t) = \cos(0) = 1
\]

对应的系数为：

\[
\dfrac{A}{\pi} \cdot \dfrac{1}{1} = \dfrac{A}{\pi}
\]

因此，直流分量为：

\[
S_{\text{DC,out}} = \dfrac{A}{\pi}
\]

锁相放大器的输出电压 \( V_{\text{out}} \) 与输入信号中与参考信号同频且同相的分量成正比：

\[
V_{\text{out}} \propto \langle S_{\text{chop}}(t) \cdot R(t) \rangle
\]

其中 \( \langle \cdot \rangle \) 表示经过低通滤波器后的平均值。

虽然在乘法器输出后出现了直流分量 \( \dfrac{A}{\pi} \)，但这个直流分量来源于高频信号的互相作用（例如 \( n = 1 \) 的谐波与参考信号相乘得到的结果）。这些直流分量与输入信号中我们关心的那部分（与参考信号同频且同相的分量）无关。

锁相放大器的核心在于相敏检测，即它只对输入信号中与参考信号频率和相位匹配的分量敏感。

我们可以计算锁相放大器输出对直流分量的响应：

锁相放大器输出：

\[
V_{\text{out}} \propto \langle S_{\text{chop}}(t) \cdot R(t) \rangle = \langle A \cdot M(t) \cdot \sin(2\pi f_{\text{chop}} t) \rangle
\]

将 \( M(t) \) 代入：

\[
V_{\text{out}} \propto A \left\langle \left( \dfrac{1}{2} + \dfrac{2}{\pi} \sum_{n=1,3,5,\dots}^{\infty} \dfrac{1}{n} \sin(2\pi n f_{\text{chop}} t) \right) \sin(2\pi f_{\text{chop}} t) \right\rangle
\]

展开后：

\[
\begin{aligned}
	V_{\text{out}} \propto & A \left\langle \dfrac{1}{2} \sin(2\pi f_{\text{chop}} t) \right\rangle \\
	& + A \left\langle \dfrac{2}{\pi} \sum_{n=1,3,5,\dots}^{\infty} \dfrac{1}{n} \sin(2\pi n f_{\text{chop}} t) \sin(2\pi f_{\text{chop}} t) \right\rangle
\end{aligned}
\]

计算平均值（低通滤波）时，任何随时间振荡的项，其平均值为零。

第一项平均值：

\[
\left\langle \dfrac{1}{2} \sin(2\pi f_{\text{chop}} t) \right\rangle = 0
\]

第二项平均值，对于每个 \( n \)：

\[
\left\langle \sin(2\pi n f_{\text{chop}} t) \sin(2\pi f_{\text{chop}} t) \right\rangle = \dfrac{1}{2} \left\langle \cos(2\pi (n - 1) f_{\text{chop}} t) - \cos(2\pi (n + 1) f_{\text{chop}} t) \right\rangle
\]

再次，由于余弦函数的平均值为零（除非频率为零的情况），因此，当 \( n = 1 \) 时：

\[
\cos(2\pi (1 - 1) f_{\text{chop}} t) = \cos(0) = 1
\]

对应的平均值为 1，但其他项的平均值均为零。因此，仅当 \( n = 1 \) 时，第二项对平均值有贡献。

计算 \( n = 1 \) 时的贡献：

\[
V_{\text{out}} \propto A \cdot \dfrac{2}{\pi} \cdot \dfrac{1}{1} \cdot \dfrac{1}{2} \cdot 1 = \dfrac{A}{\pi}
\]

这是一个直流值。虽然数学上得到一个直流分量 \( \dfrac{A}{\pi} \)，但这个值是由于信号内部高频分量相互作用的结果，而不是原始信号中与参考信号同频且同相的分量。锁相放大器的设计目的是测量输入信号中与参考信号严格同步的分量。因此，它会忽略由于其他频率分量相互作用而产生的直流偏置。

在实际测量中，直流偏置可能会引入仪器的零点漂移或偏置。然而，锁相放大器通常具有直流耦合和交流耦合模式，以及直流偏置调零功能，可以消除这些影响。锁相放大器通过同步检测，仅提取与参考信号频率和相位匹配的信号分量。其他频率（包括直流）和相位的信号分量会被滤除或平均为零。

% 思考题
\begin{question}
	相位以及相位差的含义是什么？锁相放大器输出的$\theta$是待测信号的相位还是待测信号与参考信号之间的相位差？
\end{question}
在信号处理中，相位是描述正弦波形位置的一个参数，表示一个周期性信号在时间轴上的特定点。通常一个正弦信号可以表示为：

\[
V(t) = A \sin(2\pi f t + \phi)
\]

其中，\( A \) 是信号的幅度，\( f \) 是信号的频率，\( \phi \) 是信号的相位，即信号的初始相位角，决定了信号在时间 \( t = 0 \) 时的位置。

相位 \( \phi \) 的单位是弧度或度数，它表示正弦波在时间零点相对于完整周期的位置。如果 \( \phi = 0 \)，则信号从零开始上升。如果 \( \phi = \frac{\pi}{2} \)，则信号从其最大值开始。

相位差是两个正弦波形之间的相位角的差异，它表示一个信号相对于另一个信号的时间偏移。假设有两个正弦信号：

\[
V_1(t) = A_1 \sin(2\pi f t + \phi_1)
\]
\[
V_2(t) = A_2 \sin(2\pi f t + \phi_2)
\]

那么这两个信号之间的相位差为：

\[
\Delta \phi = \phi_1 - \phi_2
\]

相位差可以正也可以负，表示一个信号相对于另一个信号是提前还是滞后。如果相位差为零，说明两个信号是同步的，波峰和波谷完全对齐；如果相位差为 \( \pi \)，则两个信号相位相反，即波峰对波谷。

锁相放大器的相位测量涉及两个重要概念。待测信号的相位：待测信号是一个周期性信号，其相位 \( \phi_{\text{signal}} \) 表示它在时间轴上的特定位置。参考信号的相位：参考信号是锁相放大器内部生成的，通常与实验中某个已知的调制信号同步。它的相位 \( \phi_{\text{ref}} \) 由系统定义，并作为相位测量的基准。

锁相放大器在测量过程中通过比较待测信号和参考信号的相位，输出两者的相位差：

\[
\Delta \phi = \phi_{\text{signal}} - \phi_{\text{ref}}
\]

锁相放大器会通过乘法器将待测信号与参考信号相乘，提取出同频率信号中的相位差信息。最终经过低通滤波器，得到一个与相位差 \( \Delta \phi \) 相关的输出。这个相位差 \( \Delta \phi \) 是锁相放大器测量的重要参数。

锁相放大器输出的 θ 通常指的是待测信号与参考信号之间的相位差。具体地说，锁相放大器通过计算待测信号和参考信号在同频率下的相位差，输出一个相位角 \( \theta \)，这个值表示待测信号相对于参考信号的偏移程度。

锁相放大器通过同时生成正交的两个参考信号，一个是与待测信号同步的正弦信号 \( \sin(2\pi f_{\text{ref}} t) \)，另一个是相位提前 \( \frac{\pi}{2} \) 的余弦信号 \( \cos(2\pi f_{\text{ref}} t) \)。然后，它分别将待测信号与这两个参考信号相乘：

\[
X = A_{\text{signal}} \cos(\phi_{\text{signal}} - \phi_{\text{ref}})
\]
\[
Y = A_{\text{signal}} \sin(\phi_{\text{signal}} - \phi_{\text{ref}})
\]

这里 \( X \) 和 \( Y \) 分别是“同相分量”和“正交分量”。

然后，锁相放大器可以根据这两个分量计算出相位差：

\[
\theta = \text{atan2}(Y, X) = \phi_{\text{signal}} - \phi_{\text{ref}}
\]

这个输出的 \( \theta \) 就是待测信号与参考信号之间的相位差。

% 思考题
\begin{question}
	锁相放大器的使用说明、操作视频和仿真软件另外给出，帮助预习。
\end{question}
详见参考文献[3]。

% 思考题
\begin{question}
	（选）以上仅讨论了滤波器的幅频特性，那么其相频特性又会是怎样呢？在时域的时间响应特性又是怎样呢？感兴趣的同学可以自己推导。进一步，该相频特性会影响到相位角的测量吗？
\end{question}
详见实验原理（Principle）部分。

\subsubsection{Reflection Questions of Preview}

% 思考题
\begin{question}
	如何用锁相放大器检测到待测的直流信号或慢变信号？
\end{question}
这个思考题在前面（Reflection Question 1.2中）已经完成了，此处不再赘述。

% 思考题
\begin{question}
	噪声有哪些类型？一般测量对象本身的噪声是哪来的？它有什么特征？
\end{question}
\begin{enumerate}
	\item 噪声的类型
	\begin{enumerate}
		\item 热噪声（Johnson噪声）：热噪声是由导体中的电子热运动引起的。根据热力学原理，当温度不为零时，导体中的电子会因热运动而产生随机的电压波动。热噪声与温度成正比，与电阻有关，且是白噪声，即在所有频率上都有相同的功率谱密度。热噪声的功率可以用公式表示：
		
		\[
		V_{\text{rms}} = \sqrt{4k_B T R B}
		\]
		
		其中，\( k_B \) 是玻尔兹曼常数，\( T \) 是温度，\( R \) 是电阻，\( B \) 是带宽。
		
		\item 1/f 噪声（闪烁噪声或低频噪声）：1/f噪声常见于半导体器件、金属导体、电阻等材料中。它与电子的迁移、材料中的缺陷和电荷俘获与释放等机制有关。1/f噪声的功率谱密度随频率增加而减小，一般表现为随频率 \( f \) 的倒数变化，即功率谱密度与 \( 1/f \) 成正比。它在低频时显著，在高频时逐渐减小，因此对低频信号测量影响较大。典型的功率谱密度可以表示为：
		 
		\[
		S(f) \propto \frac{1}{f^k}
		\]
		 
		其中 \( k \approx 1 \)。
		
		\item 散粒噪声（Shot Noise）：散粒噪声是由于电流中的电荷离散性造成的，它源自电荷载流子（如电子或空穴）在通过一个势垒时产生的离散变化。它通常出现在半导体器件、光电器件和真空管中。散粒噪声也是白噪声，它的功率谱密度与电流成正比。散粒噪声的均方根电压为：
		 
		\[
		I_{\text{shot}} = \sqrt{2qI B}
		\]
		
		其中，\( q \) 是电子电荷，\( I \) 是平均电流，\( B \) 是带宽。
		
		\item 量子噪声：量子噪声是由量子力学效应引起的，尤其是在非常低温或极高精度的测量中变得显著。比如在光电探测器中，单光子的随机到达会引起量子噪声。量子噪声通常与设备的测量极限相关，特别是在精密测量中。它的特征依赖于设备类型和所涉及的量子现象。
		
		\item 环境噪声：环境噪声来自外部设备、功率线、电磁辐射、机械振动等。最常见的例子是电源噪声，例如50Hz或60Hz的市电干扰。环境噪声通常是特定频率的噪声或周期性噪声，可以通过适当的屏蔽、接地或滤波进行消除。
	\end{enumerate}
	
	\item 测量对象本身的噪声来源
	
	测量对象本身的噪声通常源自其内在物理机制，这类噪声与测量设备无关，即它是由被测量对象的固有特性引起的。主要来源包括：
	
	在导体或半导体材料中，电子的热运动会引起电压波动。任何具有电阻的材料（如电阻器或半导体器件）都可能产生热噪声。热噪声是白噪声，其幅度与温度、阻值成正比。因此，当测量温度较高的物体或使用高阻值设备时，热噪声会变得显著。
	
	许多半导体器件和材料中的电子迁移过程、缺陷俘获和释放机制会导致1/f噪声。这种噪声广泛存在于各种电子器件中，尤其是当信号频率较低时。1/f 噪声在低频段非常明显，因此它主要影响慢变信号或直流信号的测量。
	
	在涉及电荷传输的设备中，如光电二极管或半导体器件，电流的离散性会导致散粒噪声。散粒噪声与电流的平均值成正比，通常是白噪声，影响高精度电流或光强度测量。
	
	当测量对象涉及到非常小的量子效应时，例如单光子探测、低温超导设备或精密量子系统中的测量，量子噪声会变得明显。量子噪声受系统的量子态和测量的量子不确定性限制，通常在极高精度实验中对信号的测量极限起到决定性作用。
	
	\item 测量对象本身噪声的特征
	
	噪声的本质是随机波动，这意味着它没有固定的模式，且难以完全消除。无论是热噪声、1/f噪声还是散粒噪声，它们的波动都是非周期的，频谱上表现为连续分布。
	
	热噪声和散粒噪声都是白噪声，它们在整个频谱中功率谱密度均匀分布。1/f噪声则在低频处具有更高的功率谱密度，在高频处逐渐衰减。因此，1/f噪声主要影响低频测量。量子噪声受限于测量系统的量子性质，具有更复杂的频率特征，通常与特定的量子过程相关。
	
	测量对象本身的噪声是固有的，因此无法完全消除。提高系统测量精度的关键在于减小外部噪声干扰和优化设备性能来抑制或滤除噪声。
\end{enumerate}

% 思考题
\begin{question}
	是否可以用RIGOL DG4162信号发生器产生白噪声取代教学实验箱产生的白噪声？如何将它与信号混合？
\end{question}
\begin{itemize}
	\item RIGOL DG4162 是一款多功能的信号发生器，具备产生各种波形和噪声信号的能力。它可以产生白噪声，即可以用它取代教学实验箱产生的白噪声。
	
	RIGOL DG4162 支持白噪声的输出。白噪声是一种具有恒定功率谱密度的信号，意味着它在每个频率上都具有相同的能量。根据 RIGOL DG4162 的规格说明，它可以生成多种常见的噪声信号，包括白噪声。可以通过以下步骤在 DG4162 上生成白噪声：进入信号发生菜单--选择信号类型为Noise（噪声）--调整输出幅度、偏置和频率范围等参数。
	
	\item 在实验中，常常需要将白噪声与某一待测信号（例如正弦波、方波、三角波等）进行混合，以模拟实际环境中的噪声影响。要完成这一操作，可以通过以下几种方式进行信号混合：
	\begin{enumerate}
		\item 使用外部电路混合（硬件混合）：从 RIGOL DG4162 的两个输出通道分别产生白噪声和待测信号。例如，使用 CH1 产生正弦信号，CH2 产生白噪声。通过一个运算放大器配置成加法电路，将两个信号进行混合。该加法电路可以实现输入信号的线性叠加，使得输出信号包含了两个输入信号的叠加效果。典型的电路如下，该输出信号将同时包含待测信号和白噪声：
		 
		\[
		V_{\text{out}} = V_{\text{signal}} + V_{\text{noise}}
		\]
		
		\item  信号发生器内部混合（软件混合）：
		
		一些高级信号发生器支持信号叠加功能，即可以在内部将多个信号合成，然后通过单一输出通道发送。在RIGOL DG4162中，可以通过“信号叠加”功能实现信号混合：首先选择一个通道（例如 CH1）输出待测信号，例如正弦波或方波。在主信号设置完成后，进入“噪声”菜单，选择将噪声叠加到主信号中。DG4162 可以允许设置噪声的幅度，以控制噪声与主信号的比例。
	\end{enumerate}
\end{itemize}

% 思考题
\begin{question}
	复习用示波器测量信号的实验步骤，学习用数字示波器中的带通滤波器过滤噪声的操作。
\end{question}
在实验中，噪声会干扰信号的测量，尤其是低频 $1/f$ 噪声或高频干扰噪声。数字示波器通常具备带通滤波器功能，可以有效过滤掉不需要的频段噪声。

\begin{enumerate}
	\item 什么是带通滤波器？
	
	带通滤波器是一种可以通过特定频率范围信号的滤波器，它允许在一个频段内的信号通过，同时衰减低于和高于该频段的信号。使用带通滤波器可以有效去除信号中的低频和高频噪声，只保留感兴趣的频率部分。
	
	\item 使用带通滤波器的步骤
	
	以下是通过数字示波器使用带通滤波器的基本步骤：
	
	\begin{enumerate}
		\item 确定感兴趣的频段
		
		首先确定想观测或测量的信号的频率范围。例如，如果测量的信号是 $1$kHz 的正弦波，那么可以设置带通滤波器的通带范围为 $500$Hz 至 $2$kHz，这样可以过滤掉低频和高频的噪声干扰。
		
		\item 进入滤波器设置界面
		\begin{itemize}
			\item 在数字示波器中，通常会有“滤波”选项。在 RIGOL 等品牌的示波器中，可以通过菜单选择信号处理（Signal Processing）或数学操作（Math Operations），然后选择“带通滤波器”（Band-Pass Filter）选项。
		\end{itemize}
		
		\item 设置滤波器参数
		\begin{itemize}
			\item 设置带通滤波器的 \textbf{下截止频率} 和 \textbf{上截止频率}。例如，如果想保留 $1$kHz 附近的信号并消除其他频段的噪声，设置下截止频率为 $500$Hz，上截止频率为 $2$kHz。
			\item 可以根据实际测量需求精细调整这些截止频率，以获得最佳的信号净化效果。
		\end{itemize}
		
		\item 启用滤波器
		
		确认滤波器设置后，启用带通滤波器。数字示波器会自动应用滤波操作，将看到屏幕上显示的信号减少了噪声干扰，显示更加平滑和清晰。
	\end{enumerate}
	
	\item 滤波后信号的分析
	\begin{itemize}
		\item 在带通滤波器应用后，观察滤波后的信号。可以再次使用示波器的测量功能，检查信号的频率、幅度和相位等参数是否符合预期。
		\item 带通滤波器可以有效过滤不必要的噪声，但需要注意滤波器可能会对信号造成一些相位延迟，尤其是在信号的边缘部分。因此，如果相位信息对测量结果至关重要，需谨慎选择滤波参数。
	\end{itemize}
\end{enumerate}

